\documentclass{situationist-preprint}

\title{Hilbert's Hotel}
\msid{MTH-10-b}
\msfamily{MTH-10}
\msdiscipline{Philosophy of Mathematics}
\mstype{Companion manuscript}
\mskeywords{actual infinity; Hilbert's hotel; part--whole principle; kalām; transfinite arithmetic}

\begin{document}
\maketitle

\begin{abstract}
Hilbert's hotel is offered as a reductio of actual infinity. We examine it in detail: each scenario describes a bijection between $\omega$ and a part of it, provable in set theory. The hotel violates the part--whole principle, which fails for infinite collections, not any law of logic. This supports MTH-10-a's conclusion that actual infinities exist.
\end{abstract}

\paragraph{The scenarios.} A hotel has a room for each natural number, all occupied. (i) A new guest arrives: each guest moves from room $n$ to room $n+1$, and room $0$ is free. (ii) Countably many guests arrive: each moves from $n$ to $2n$, and the odd rooms are free. (iii) Countably many coaches arrive, each with countably many passengers: guest $n$ moves to $2^{n}$, and passenger $m$ of coach $k$ takes room $p_k^{m+1}$, where $p_k$ is the $(k+1)$th odd prime. Each scenario is a one-one map of $\omega$ into $\omega$ that is not onto, provably existing in set theory. The story, popularized by \citet{Gamow1947}, exhibits no contradiction.

\paragraph{The diagnosis.} The air of absurdity comes from the principle that a whole is greater than its proper part, which holds for finite collections and fails, by Dedekind's definition, for infinite ones, as MTH-10-a (Section~4) argues. \citet{Mancosu2009} shows that sizes respecting the part--whole principle can be consistently defined for sets of natural numbers, but not together with the equation of size with one-one correspondence. Neither choice yields a contradiction.

\begin{objection}
\citet{Craig1979} adds departures: if the guests in the odd rooms leave, infinitely many remain; if all from room $3$ onward leave, three remain. Identical subtractions give different results.
\end{objection}

\begin{reply}
Cardinal subtraction is undefined for infinite cardinals, as division by zero is; each departure removes a different set and leaves a determinate remainder. An undefined operation is not a contradiction.
\end{reply}

\paragraph{Conclusion.} Hilbert's hotel is a model of $\omega$, not a reductio of it.

\begin{thebibliography}{99}\small
\bibitem[Craig(1979)]{Craig1979} Craig, W.~L. (1979). \emph{The Kalām Cosmological Argument}. Macmillan.
\bibitem[Gamow(1947)]{Gamow1947} Gamow, G. (1947). \emph{One Two Three\ldots Infinity}. Viking Press.
\bibitem[Mancosu(2009)]{Mancosu2009} Mancosu, P. (2009). Measuring the size of infinite collections of natural numbers: Was Cantor's theory of infinite number inevitable? \emph{The Review of Symbolic Logic}, 2(4), 612--646.
\end{thebibliography}

\end{document}
